% Figure IV.10 --- the three-step settlement protocol S:fh-settle-proto
% describes (bond post, damage claim, verify-and-clearance), drawn as a
% three-lifeline sequence with escrow as its own lifeline -- distinct from
% fig:fh-settlement, which draws the custody-bound counterfactual
% (S:fh-escrow), not this protocol. Fixes the corpus defect where
% fig-fh-settlement.tex was pressed into service for two different claims
% under one label; this fragment gives the settlement-protocol prose its own
% figure and its own label.
%
% Reader question: which of the three settlement steps can a signed
%   inclusion proof alone not settle by itself, and why?
% Claim: the protocol is exactly three steps in strict order -- bond post,
%   damage claim, verify-and-clearance -- and only the third step is a
%   decision the escrow makes; the first two are one-sided announcements.
% Evidence: the chapter's own three-step enumeration in S:fh-settle-proto:
%   the bond B_alpha locked at post time, the claim
%   cl_B=(R_B^(e),evidence), and the two terminal outcomes clear/refuse from
%   Design Invariant thm:fh-escrow-bound.
% Must distinguish: the two one-sided announcements (post, claim) from the
%   one step where the escrow actually decides (verify), and that the
%   decision has exactly two terminal outcomes, not a spectrum.
% Grammar chosen: three-lifeline sequence (A, escrow, B), reusing
%   fig-fh-settlement's own Clear/Refuse fan for the terminal step so the
%   two figures read as one family without redrawing the same schema.
%   Grammar rejected: a three-row table would lose the fact that step 3
%   branches into two outcomes while 1--2 do not; a state machine would
%   overstate this as a machine with memory beyond three announcements plus
%   one decision.
% Five-second test: which step is the only one with two possible outcomes?
\pdfigurehue{pdgold}%
\begin{figure}[H]
\centering
\begin{tikzpicture}[pd figure,x=1cm,y=1cm]
  \def\xA{0.90}  \def\xG{5.20}  \def\xB{9.50}
  \node[pd title,anchor=west,text width=9.0cm] at (0.00,1.00)
    {two announcements, then one decision};
  \node[pd state,minimum width=2.6cm] (hA) at (\xA,0) {Harbor $A$};
  \node[draw=pdinkmuted!70,line width=.9pt,fill=pdpage,diamond,
    minimum width=1.5cm,minimum height=0.85cm,inner sep=0pt,pd label]
    (hG) at (\xG,0) {escrow};
  \node[pd state,minimum width=2.6cm] (hB) at (\xB,0) {Harbor $B$};
  \draw[pd guide] (\xA,-0.60) -- (\xA,-4.55);
  \draw[pd guide] (\xG,-0.60) -- (\xG,-4.15);
  \draw[pd guide] (\xB,-0.60) -- (\xB,-4.55);

  \node[pd tag,anchor=south,text width=4.4cm,align=center] at ({(\xA+\xG)/2},-1.16)
    {bond post: $B_\alpha$ locked, signed into $R_A{}^{(e)}$};
  \node[pd badge] (m1) at ({(\xA+\xG)/2},-1.40) {1};
  \draw[pd arrow,pd to badge] (\xA,-1.20) -- (m1);
  \draw[pd arrow,pd from badge] (m1) -- (\xG,-1.60);

  \node[pd tag,anchor=south,text width=4.4cm,align=center] at ({(\xG+\xB)/2},-2.36)
    {damage claim: $\mathsf{cl}_B=(R_B{}^{(e)},\mathrm{evidence})$};
  \node[pd badge] (m2) at ({(\xG+\xB)/2},-2.60) {2};
  \draw[pd arrow,pd to badge] (\xB,-2.40) -- (m2);
  \draw[pd arrow,pd from badge] (m2) -- (\xG,-2.80);

  \node[draw=pdinkmuted!70,line width=.9pt,fill=pdpage,diamond,
    minimum width=1.6cm,minimum height=0.90cm,inner sep=0pt,pd label]
    (verify) at (\xG,-3.70) {verify};
  \node[pd badge] at (\xG-1.35,-3.70) {3};
  \node[pd focus datum] (clear) at (\xB-0.60,-4.90) {};
  \draw[pd focus arrow] (verify) -- (clear);
  \node[pd focus label,anchor=north,text width=4.0cm,align=left] at (\xB-1.55,-5.12)
    {\textbf{clear}: pay $B$, refund remainder to $A$};
  \node[pd datum] (refuse) at (\xA+0.60,-4.90) {};
  \draw[pd arrow] (verify) -- (refuse);
  \node[pd label,anchor=north,text width=3.6cm,align=right] at (\xA+1.55,-5.12)
    {\textbf{refuse}: return bond to $A$};
\end{tikzpicture}

\vspace{4pt}
\footnotesize
\renewcommand{\arraystretch}{1.12}
\begin{tabular}{@{}c l p{6.3cm}@{}}
\toprule
\# & step & binds / value \\
\midrule
1 & bond post & $A\to$ escrow: $B_\alpha$ locked, TTL signed into $R_A{}^{(e)}$ and $R_B{}^{(e)}$ \\
2 & damage claim & $B\to$ escrow: $\mathsf{cl}_B=(R_B{}^{(e)},\mathrm{evidence})$, a Merkle inclusion proof \\
3 & verify \& clear & escrow checks the proof against step 1's terms; clear or refuse \\
\bottomrule
\end{tabular}
\normalsize
\caption{Posting a bond and announcing damage do not settle it. Escrow checks the proof against the signed conditions, then makes one terminal decision: clear or refuse. [internal]}
\label{fig:fh-settle-protocol}
\end{figure}
